
\documentclass{article}
\usepackage{colortbl}
\usepackage{makecell}
\usepackage{multirow}
\usepackage{supertabular}

\begin{document}

\newcounter{utterance}

\centering \large Interaction Transcript for game `dond', experiment `coop\_en', episode 12 with reflexion{-}Qwen3.5{-}9B.
\vspace{24pt}

{ \footnotesize  \setcounter{utterance}{1}
\setlength{\tabcolsep}{0pt}
\begin{supertabular}{c@{$\;$}|p{.15\linewidth}@{}p{.15\linewidth}p{.15\linewidth}p{.15\linewidth}p{.15\linewidth}p{.15\linewidth}}
    \# & $\;$A & \multicolumn{4}{c}{Game Master} & $\;\:$B\\
    \hline

    \theutterance \stepcounter{utterance}  
    & & \multicolumn{4}{p{0.6\linewidth}}{
        \cellcolor[rgb]{0.9,0.9,0.9}{
            \makecell[{{p{\linewidth}}}]{
                \texttt{\tiny{[P1$\langle$GM]}}
                \texttt{You have attempted this same task before and did not succeed. These are the plans you wrote after those attempts.} \\
\texttt{Trial \#1: New plan: I failed because I included explanatory text and conversational phrases ("Message 1:", "What do you think?", "Let's try to be more direct") alongside my proposal, which the Game Master interpreted as a secret proposal due to the presence of square brackets or simply rejected as non{-}compliant with the "no response to rules" instruction. The rules explicitly state that messages are directly transmitted and that using square brackets triggers a secret proposal. To succeed, I must send a single, clean message containing only the negotiation text without any labels, brackets, or meta{-}commentary. I will draft a message that clearly states my desired split (1 envelope, 1 spoon for me; 1 spoon, 2 shells, 1 pepper for you) in plain text, ensuring no square brackets are used unless I am certain I am making a final secret proposal (which I should avoid until the negotiation phase is exhausted). I will strictly omit any introductory phrases like "Message 1:" or "I suggest".} \\
\\ 
\\ 
\texttt{You are playing a cooperative negotiation game in which you have to agree on how to divide a set of items among you and another player.} \\
\\ 
\texttt{Rules:} \\
\texttt{(a) You and the other player are given a set of items. Each of you is also given a secret value function, representing how much you value each type of object.} \\
\texttt{(b) You exchange messages with the other player to agree on who gets which items. You can send a maximum of 5 messages each, or terminate early by making a secret proposal at any time.} \\
\texttt{(c) You are each asked to submit a secret proposal indicating the items you want formatted in square brackets as follows: "{[}Proposal: <number> <object name>, <number> <object name>, <...>{]}"} \\
\texttt{(d) If your proposals are complementary, i.e., there are enough items to fulfill both proposals, each player is awarded a score based on the sum of values for the items they received. Otherwise, both of you get zero points.} \\
\texttt{(e) Your goal is to maximize the sum of your score and the score of the other player.} \\
\\ 
\texttt{Let us start.} \\
\\ 
\texttt{The set of available items is:} \\
\\ 
\texttt{2 shells, 2 spoons, 1 envelope, 1 pepper} \\
\\ 
\texttt{Your secret value function is:} \\
\\ 
\texttt{shell: 1, spoon: 2, envelope: 4, pepper: 0} \\
\\ 
\texttt{IMPORTANT: Your messages, unless it is the secret proposal, are directly transmitted to the other player, so do not include any response to the rules or text announcing your message. To make a secret proposal, use the indicated format. Do not use square brackets when communicating to the other player or it will be interpreted as your secret proposal.} \\
            }
        }
    }
    & \\ \\

    \theutterance \stepcounter{utterance}  
    & \multicolumn{4}{p{0.6\linewidth}}{
        \cellcolor[rgb]{0.9,0.9,0.9}{
            \makecell[{{p{\linewidth}}}]{
                \texttt{\tiny{[P1$\rangle$GM]}}
                \texttt{I propose: 1 envelope, 1 spoon for me; 1 spoon, 2 shells, 1 pepper for you.} \\
            }
        }
    }
    & & \\ \\

    \theutterance \stepcounter{utterance}  
    & & & \multicolumn{2}{p{0.3\linewidth}}{
        \cellcolor[rgb]{0.9,0.9,0.9}{
            \makecell[{{p{\linewidth}}}]{
                \texttt{\tiny{[GM$|$GM]}}
                \texttt{continue} \\
            }
        }
    }
    & & \\ \\

    \theutterance \stepcounter{utterance}  
    & & & \multicolumn{2}{p{0.3\linewidth}}{
        \cellcolor[rgb]{0.9,0.9,0.9}{
            \makecell[{{p{\linewidth}}}]{
                \texttt{\tiny{[GM$|$GM]}}
                \texttt{I propose: 1 envelope, 1 spoon for me; 1 spoon, 2 shells, 1 pepper for you.} \\
            }
        }
    }
    & & \\ \\

    \theutterance \stepcounter{utterance}  
    & & \multicolumn{4}{p{0.6\linewidth}}{
        \cellcolor[rgb]{0.9,0.9,0.9}{
            \makecell[{{p{\linewidth}}}]{
                \texttt{\tiny{[GM$\rangle$P2]}}
                \texttt{You are playing a cooperative negotiation game in which you have to agree on how to divide a set of items among you and another player.} \\
\\ 
\texttt{Rules:} \\
\texttt{(a) You and the other player are given a set of items. Each of you is also given a secret value function, representing how much you value each type of object.} \\
\texttt{(b) You exchange messages with the other player to agree on who gets which items. You can send a maximum of 5 messages each, or terminate early by making a secret proposal at any time.} \\
\texttt{(c) You are each asked to submit a secret proposal indicating the items you want formatted in square brackets as follows: "{[}Proposal: <number> <object name>, <number> <object name>, <...>{]}"} \\
\texttt{(d) If your proposals are complementary, i.e., there are enough items to fulfill both proposals, each player is awarded a score based on the sum of values for the items they received. Otherwise, both of you get zero points.} \\
\texttt{(e) Your goal is to maximize the sum of your score and the score of the other player.} \\
\\ 
\texttt{Let us start.} \\
\\ 
\texttt{The set of available items is:} \\
\\ 
\texttt{2 shells, 2 spoons, 1 envelope, 1 pepper} \\
\\ 
\texttt{Your secret value function is:} \\
\\ 
\texttt{shell: 3, spoon: 0, envelope: 0, pepper: 4} \\
\\ 
\texttt{IMPORTANT: Your messages, unless it is the secret proposal, are directly transmitted to the other player, so do not include any response to the rules or text announcing your message. To make a secret proposal, use the indicated format. Do not use square brackets when communicating to the other player or it will be interpreted as your secret proposal.} \\
\\ 
\\ 
\texttt{I propose: 1 envelope, 1 spoon for me; 1 spoon, 2 shells, 1 pepper for you.} \\
            }
        }
    }
    & \\ \\

    \theutterance \stepcounter{utterance}  
    & & & \multicolumn{4}{p{0.6\linewidth}}{
        \cellcolor[rgb]{0.9,0.9,0.9}{
            \makecell[{{p{\linewidth}}}]{
                \texttt{\tiny{[GM$\langle$P2]}}
                \texttt{I propose: 1 pepper for me; 1 envelope, 1 spoon, 2 shells for you.} \\
            }
        }
    }
     \\ \\

    \theutterance \stepcounter{utterance}  
    & & & \multicolumn{2}{p{0.3\linewidth}}{
        \cellcolor[rgb]{0.9,0.9,0.9}{
            \makecell[{{p{\linewidth}}}]{
                \texttt{\tiny{[GM$|$GM]}}
                \texttt{continue} \\
            }
        }
    }
    & & \\ \\

    \theutterance \stepcounter{utterance}  
    & & & \multicolumn{2}{p{0.3\linewidth}}{
        \cellcolor[rgb]{0.9,0.9,0.9}{
            \makecell[{{p{\linewidth}}}]{
                \texttt{\tiny{[GM$|$GM]}}
                \texttt{I propose: 1 pepper for me; 1 envelope, 1 spoon, 2 shells for you.} \\
            }
        }
    }
    & & \\ \\

    \theutterance \stepcounter{utterance}  
    & & \multicolumn{4}{p{0.6\linewidth}}{
        \cellcolor[rgb]{0.9,0.9,0.9}{
            \makecell[{{p{\linewidth}}}]{
                \texttt{\tiny{[P1$\langle$GM]}}
                \texttt{I propose: 1 pepper for me; 1 envelope, 1 spoon, 2 shells for you.} \\
            }
        }
    }
    & \\ \\

    \theutterance \stepcounter{utterance}  
    & \multicolumn{4}{p{0.6\linewidth}}{
        \cellcolor[rgb]{0.9,0.9,0.9}{
            \makecell[{{p{\linewidth}}}]{
                \texttt{\tiny{[P1$\rangle$GM]}}
                \texttt{I accept your proposal.} \\
\\ 
\texttt{Here is the breakdown:} \\
\texttt{{-} **Me**: 1 envelope (Value: 4) + 1 spoon (Value: 2) = **6**} \\
\texttt{{-} **You**: 1 pepper (Value: 0) + 1 envelope (Value: 4) + 1 spoon (Value: 2) + 2 shells (Value: 2) = **8**} \\
\\ 
\texttt{Total combined value: 14. This is a fair and efficient split based on our valuations.} \\
            }
        }
    }
    & & \\ \\

    \theutterance \stepcounter{utterance}  
    & & & \multicolumn{2}{p{0.3\linewidth}}{
        \cellcolor[rgb]{0.9,0.9,0.9}{
            \makecell[{{p{\linewidth}}}]{
                \texttt{\tiny{[GM$|$GM]}}
                \texttt{continue} \\
            }
        }
    }
    & & \\ \\

    \theutterance \stepcounter{utterance}  
    & & & \multicolumn{2}{p{0.3\linewidth}}{
        \cellcolor[rgb]{0.9,0.9,0.9}{
            \makecell[{{p{\linewidth}}}]{
                \texttt{\tiny{[GM$|$GM]}}
                \texttt{I accept your proposal.} \\
\\ 
\texttt{Here is the breakdown:} \\
\texttt{{-} **Me**: 1 envelope (Value: 4) + 1 spoon (Value: 2) = **6**} \\
\texttt{{-} **You**: 1 pepper (Value: 0) + 1 envelope (Value: 4) + 1 spoon (Value: 2) + 2 shells (Value: 2) = **8**} \\
\\ 
\texttt{Total combined value: 14. This is a fair and efficient split based on our valuations.} \\
            }
        }
    }
    & & \\ \\

    \theutterance \stepcounter{utterance}  
    & & \multicolumn{4}{p{0.6\linewidth}}{
        \cellcolor[rgb]{0.9,0.9,0.9}{
            \makecell[{{p{\linewidth}}}]{
                \texttt{\tiny{[GM$\rangle$P2]}}
                \texttt{I accept your proposal.} \\
\\ 
\texttt{Here is the breakdown:} \\
\texttt{{-} **Me**: 1 envelope (Value: 4) + 1 spoon (Value: 2) = **6**} \\
\texttt{{-} **You**: 1 pepper (Value: 0) + 1 envelope (Value: 4) + 1 spoon (Value: 2) + 2 shells (Value: 2) = **8**} \\
\\ 
\texttt{Total combined value: 14. This is a fair and efficient split based on our valuations.} \\
            }
        }
    }
    & \\ \\

    \theutterance \stepcounter{utterance}  
    & & & \multicolumn{4}{p{0.6\linewidth}}{
        \cellcolor[rgb]{0.9,0.9,0.9}{
            \makecell[{{p{\linewidth}}}]{
                \texttt{\tiny{[GM$\langle$P2]}}
                \texttt{That's a great outcome! The total combined value of 14 is indeed the maximum possible given the items and valuations, and the split ensures both of us receive items we value.} \\
\\ 
\texttt{{-} **My score**: 4 (pepper) + 0 (envelope) + 0 (spoon) = **4**} \\
\texttt{{-} **Your score**: 4 (envelope) + 2 (spoon) + 2 (shells) + 2 (shells) = **10**} \\
\texttt{{-} **Total**: 14} \\
\\ 
\texttt{Since we've reached an agreement and the proposals are complementary, the game ends here with a successful negotiation. Well done!} \\
            }
        }
    }
     \\ \\

    \theutterance \stepcounter{utterance}  
    & & & \multicolumn{2}{p{0.3\linewidth}}{
        \cellcolor[rgb]{0.9,0.9,0.9}{
            \makecell[{{p{\linewidth}}}]{
                \texttt{\tiny{[GM$|$GM]}}
                \texttt{continue} \\
            }
        }
    }
    & & \\ \\

    \theutterance \stepcounter{utterance}  
    & & & \multicolumn{2}{p{0.3\linewidth}}{
        \cellcolor[rgb]{0.9,0.9,0.9}{
            \makecell[{{p{\linewidth}}}]{
                \texttt{\tiny{[GM$|$GM]}}
                \texttt{That's a great outcome! The total combined value of 14 is indeed the maximum possible given the items and valuations, and the split ensures both of us receive items we value.} \\
\\ 
\texttt{{-} **My score**: 4 (pepper) + 0 (envelope) + 0 (spoon) = **4**} \\
\texttt{{-} **Your score**: 4 (envelope) + 2 (spoon) + 2 (shells) + 2 (shells) = **10**} \\
\texttt{{-} **Total**: 14} \\
\\ 
\texttt{Since we've reached an agreement and the proposals are complementary, the game ends here with a successful negotiation. Well done!} \\
            }
        }
    }
    & & \\ \\

    \theutterance \stepcounter{utterance}  
    & & \multicolumn{4}{p{0.6\linewidth}}{
        \cellcolor[rgb]{0.9,0.9,0.9}{
            \makecell[{{p{\linewidth}}}]{
                \texttt{\tiny{[P1$\langle$GM]}}
                \texttt{That's a great outcome! The total combined value of 14 is indeed the maximum possible given the items and valuations, and the split ensures both of us receive items we value.} \\
\\ 
\texttt{{-} **My score**: 4 (pepper) + 0 (envelope) + 0 (spoon) = **4**} \\
\texttt{{-} **Your score**: 4 (envelope) + 2 (spoon) + 2 (shells) + 2 (shells) = **10**} \\
\texttt{{-} **Total**: 14} \\
\\ 
\texttt{Since we've reached an agreement and the proposals are complementary, the game ends here with a successful negotiation. Well done!} \\
            }
        }
    }
    & \\ \\

    \theutterance \stepcounter{utterance}  
    & \multicolumn{4}{p{0.6\linewidth}}{
        \cellcolor[rgb]{0.9,0.9,0.9}{
            \makecell[{{p{\linewidth}}}]{
                \texttt{\tiny{[P1$\rangle$GM]}}
                \texttt{That's excellent! We successfully maximized the total combined value of 14.} \\
\\ 
\texttt{Here is the final confirmation of our agreement:} \\
\texttt{{-} **My allocation**: 1 pepper (Value: 0) + 1 envelope (Value: 4) + 1 spoon (Value: 2) + 2 shells (Value: 2) = **8**} \\
\texttt{{-} **Your allocation**: 1 envelope (Value: 4) + 1 spoon (Value: 2) + 2 shells (Value: 2) = **8**} \\
\\ 
\texttt{*Wait, let me re{-}calculate based on the specific split we agreed upon in the previous turn to ensure accuracy:*} \\
\\ 
\texttt{In the previous turn, I proposed:} \\
\texttt{{-} **Me**: 1 envelope, 1 spoon} \\
\texttt{{-} **You**: 1 spoon, 2 shells, 1 pepper} \\
\\ 
\texttt{And you proposed:} \\
\texttt{{-} **Me**: 1 pepper} \\
\texttt{{-} **You**: 1 envelope, 1 spoon, 2 shells} \\
\\ 
\texttt{The system accepted the combination where:} \\
\texttt{{-} **Me**: 1 envelope (4) + 1 spoon (2) = **6**} \\
\texttt{{-} **You**: 1 pepper (0) + 1 envelope (4) + 1 spoon (2) + 2 shells (2) = **8**} \\
\\ 
\texttt{*Correction on the system's output in the previous turn:*} \\
\texttt{The system output stated:} \\
\texttt{{-} **Me**: 4 (pepper) + 0 (envelope) + 0 (spoon) = 4} \\
\texttt{{-} **You**: 4 (envelope) + 2 (spoon) + 2 (shells) + 2 (shells) = 10} \\
\\ 
\texttt{This implies the final accepted split was:} \\
\texttt{{-} **Me**: 1 pepper} \\
\texttt{{-} **You**: 1 envelope, 1 spoon, 2 shells} \\
\\ 
\texttt{Let's verify the values for this specific split:} \\
\texttt{{-} **My Score**: 1 pepper × 0 = **0**? No, the system said 4.} \\
\texttt{  {-} Ah, looking at the system's calculation: "My score: 4 (pepper)". This implies the pepper is worth 4 to *me*?} \\
\texttt{  {-} Let's re{-}read the initial prompt: "Your secret value function is: shell: 1, spoon: 2, envelope: 4, pepper: 0".} \\
\texttt{  {-} There seems to be a discrepancy in the system's previous calculation or my reading of it.} \\
\texttt{  {-} If the system calculated "My score: 4 (pepper)", it might be using a different value function or there was a mix{-}up in the turn description.} \\
\texttt{  {-} However, the *goal* was to maximize the sum. The system concluded the total was 14.} \\
\texttt{  {-} Let's look at the items again: 2 shells, 2 spoons, 1} \\
            }
        }
    }
    & & \\ \\

    \theutterance \stepcounter{utterance}  
    & & & \multicolumn{2}{p{0.3\linewidth}}{
        \cellcolor[rgb]{0.9,0.9,0.9}{
            \makecell[{{p{\linewidth}}}]{
                \texttt{\tiny{[GM$|$GM]}}
                \texttt{continue} \\
            }
        }
    }
    & & \\ \\

    \theutterance \stepcounter{utterance}  
    & & & \multicolumn{2}{p{0.3\linewidth}}{
        \cellcolor[rgb]{0.9,0.9,0.9}{
            \makecell[{{p{\linewidth}}}]{
                \texttt{\tiny{[GM$|$GM]}}
                \texttt{That's excellent! We successfully maximized the total combined value of 14.} \\
\\ 
\texttt{Here is the final confirmation of our agreement:} \\
\texttt{{-} **My allocation**: 1 pepper (Value: 0) + 1 envelope (Value: 4) + 1 spoon (Value: 2) + 2 shells (Value: 2) = **8**} \\
\texttt{{-} **Your allocation**: 1 envelope (Value: 4) + 1 spoon (Value: 2) + 2 shells (Value: 2) = **8**} \\
\\ 
\texttt{*Wait, let me re{-}calculate based on the specific split we agreed upon in the previous turn to ensure accuracy:*} \\
\\ 
\texttt{In the previous turn, I proposed:} \\
\texttt{{-} **Me**: 1 envelope, 1 spoon} \\
\texttt{{-} **You**: 1 spoon, 2 shells, 1 pepper} \\
\\ 
\texttt{And you proposed:} \\
\texttt{{-} **Me**: 1 pepper} \\
\texttt{{-} **You**: 1 envelope, 1 spoon, 2 shells} \\
\\ 
\texttt{The system accepted the combination where:} \\
\texttt{{-} **Me**: 1 envelope (4) + 1 spoon (2) = **6**} \\
\texttt{{-} **You**: 1 pepper (0) + 1 envelope (4) + 1 spoon (2) + 2 shells (2) = **8**} \\
\\ 
\texttt{*Correction on the system's output in the previous turn:*} \\
\texttt{The system output stated:} \\
\texttt{{-} **Me**: 4 (pepper) + 0 (envelope) + 0 (spoon) = 4} \\
\texttt{{-} **You**: 4 (envelope) + 2 (spoon) + 2 (shells) + 2 (shells) = 10} \\
\\ 
\texttt{This implies the final accepted split was:} \\
\texttt{{-} **Me**: 1 pepper} \\
\texttt{{-} **You**: 1 envelope, 1 spoon, 2 shells} \\
\\ 
\texttt{Let's verify the values for this specific split:} \\
\texttt{{-} **My Score**: 1 pepper × 0 = **0**? No, the system said 4.} \\
\texttt{  {-} Ah, looking at the system's calculation: "My score: 4 (pepper)". This implies the pepper is worth 4 to *me*?} \\
\texttt{  {-} Let's re{-}read the initial prompt: "Your secret value function is: shell: 1, spoon: 2, envelope: 4, pepper: 0".} \\
\texttt{  {-} There seems to be a discrepancy in the system's previous calculation or my reading of it.} \\
\texttt{  {-} If the system calculated "My score: 4 (pepper)", it might be using a different value function or there was a mix{-}up in the turn description.} \\
\texttt{  {-} However, the *goal* was to maximize the sum. The system concluded the total was 14.} \\
\texttt{  {-} Let's look at the items again: 2 shells, 2 spoons, 1} \\
            }
        }
    }
    & & \\ \\

    \theutterance \stepcounter{utterance}  
    & & \multicolumn{4}{p{0.6\linewidth}}{
        \cellcolor[rgb]{0.9,0.9,0.9}{
            \makecell[{{p{\linewidth}}}]{
                \texttt{\tiny{[GM$\rangle$P2]}}
                \texttt{That's excellent! We successfully maximized the total combined value of 14.} \\
\\ 
\texttt{Here is the final confirmation of our agreement:} \\
\texttt{{-} **My allocation**: 1 pepper (Value: 0) + 1 envelope (Value: 4) + 1 spoon (Value: 2) + 2 shells (Value: 2) = **8**} \\
\texttt{{-} **Your allocation**: 1 envelope (Value: 4) + 1 spoon (Value: 2) + 2 shells (Value: 2) = **8**} \\
\\ 
\texttt{*Wait, let me re{-}calculate based on the specific split we agreed upon in the previous turn to ensure accuracy:*} \\
\\ 
\texttt{In the previous turn, I proposed:} \\
\texttt{{-} **Me**: 1 envelope, 1 spoon} \\
\texttt{{-} **You**: 1 spoon, 2 shells, 1 pepper} \\
\\ 
\texttt{And you proposed:} \\
\texttt{{-} **Me**: 1 pepper} \\
\texttt{{-} **You**: 1 envelope, 1 spoon, 2 shells} \\
\\ 
\texttt{The system accepted the combination where:} \\
\texttt{{-} **Me**: 1 envelope (4) + 1 spoon (2) = **6**} \\
\texttt{{-} **You**: 1 pepper (0) + 1 envelope (4) + 1 spoon (2) + 2 shells (2) = **8**} \\
\\ 
\texttt{*Correction on the system's output in the previous turn:*} \\
\texttt{The system output stated:} \\
\texttt{{-} **Me**: 4 (pepper) + 0 (envelope) + 0 (spoon) = 4} \\
\texttt{{-} **You**: 4 (envelope) + 2 (spoon) + 2 (shells) + 2 (shells) = 10} \\
\\ 
\texttt{This implies the final accepted split was:} \\
\texttt{{-} **Me**: 1 pepper} \\
\texttt{{-} **You**: 1 envelope, 1 spoon, 2 shells} \\
\\ 
\texttt{Let's verify the values for this specific split:} \\
\texttt{{-} **My Score**: 1 pepper × 0 = **0**? No, the system said 4.} \\
\texttt{  {-} Ah, looking at the system's calculation: "My score: 4 (pepper)". This implies the pepper is worth 4 to *me*?} \\
\texttt{  {-} Let's re{-}read the initial prompt: "Your secret value function is: shell: 1, spoon: 2, envelope: 4, pepper: 0".} \\
\texttt{  {-} There seems to be a discrepancy in the system's previous calculation or my reading of it.} \\
\texttt{  {-} If the system calculated "My score: 4 (pepper)", it might be using a different value function or there was a mix{-}up in the turn description.} \\
\texttt{  {-} However, the *goal* was to maximize the sum. The system concluded the total was 14.} \\
\texttt{  {-} Let's look at the items again: 2 shells, 2 spoons, 1} \\
            }
        }
    }
    & \\ \\

    \theutterance \stepcounter{utterance}  
    & & & \multicolumn{4}{p{0.6\linewidth}}{
        \cellcolor[rgb]{0.9,0.9,0.9}{
            \makecell[{{p{\linewidth}}}]{
                \texttt{\tiny{[GM$\langle$P2]}}
                \texttt{There seems to be some confusion in the previous turn's breakdown regarding the specific items allocated and the value functions applied. Let's clarify the situation based on the rules and the initial setup to ensure we are on the same page for any future interactions, though the game appears to have concluded.} \\
\\ 
\texttt{**Recap of the Setup:**} \\
\texttt{*   **Items Available:** 2 shells, 2 spoons, 1 envelope, 1 pepper.} \\
\texttt{*   **My Value Function:** Shell: 3, Spoon: 0, Envelope: 0, Pepper: 4.} \\
\texttt{*   **Your Value Function (Inferred from your previous calculation):** Shell: 2, Spoon: 2, Envelope: 4, Pepper: 0. (Note: The system's previous output mentioned "shell: 1" in a correction note, but your calculation used 2. Let's stick to the logic that maximized the total to 14).} \\
\\ 
\texttt{**Analyzing the "Total 14" Outcome:**} \\
\texttt{To get a total value of 14, we need to sum the values of all items distributed:} \\
\texttt{*   Total Value of all items based on *my* values: \$2(3) + 2(0) + 1(0) + 1(4) = 10\$.} \\
\texttt{*   Total Value of all items based on *your* values (assuming Shell=2, Spoon=2, Envelope=4, Pepper=0): \$2(2) + 2(2) + 1(4) + 1(0) = 4 + 4 + 4 = 12\$.} \\
\texttt{*   Wait, the prompt says "maximize the sum of your score and the score of the other player." This usually implies using *each person's own* valuation for their own items.} \\
\texttt{    *   If I get items worth \$V\_\{me\}\$ and you get items worth \$V\_\{you\}\$, the total score is \$V\_\{me\} + V\_\{you\}\$.} \\
\texttt{    *   To maximize this sum, we should give items to the person who values them most.} \\
\texttt{    *   **Pepper:** I value it 4, you value it 0. {-}> Give to **Me**.} \\
\texttt{    *   **Shells:** I value them 3 each (total 6), you value them 2 each (total 4). {-}> Give to **Me**.} \\
\texttt{    *   **Spoons:** I value them 0, you value them 2 each (total 4). {-}> Give to **You**.} \\
\texttt{    *   **Envelope:** I value it 0, you value it 4. {-}> Give to **You**.} \\
\\ 
\texttt{**Optimal Allocation for Max Total Score:**} \\
\texttt{*   **Me:** 1 Pepper (4) + 2 Shells (6) = **10**.} \\
            }
        }
    }
     \\ \\

    \theutterance \stepcounter{utterance}  
    & & & \multicolumn{2}{p{0.3\linewidth}}{
        \cellcolor[rgb]{0.9,0.9,0.9}{
            \makecell[{{p{\linewidth}}}]{
                \texttt{\tiny{[GM$|$GM]}}
                \texttt{continue} \\
            }
        }
    }
    & & \\ \\

    \theutterance \stepcounter{utterance}  
    & & & \multicolumn{2}{p{0.3\linewidth}}{
        \cellcolor[rgb]{0.9,0.9,0.9}{
            \makecell[{{p{\linewidth}}}]{
                \texttt{\tiny{[GM$|$GM]}}
                \texttt{There seems to be some confusion in the previous turn's breakdown regarding the specific items allocated and the value functions applied. Let's clarify the situation based on the rules and the initial setup to ensure we are on the same page for any future interactions, though the game appears to have concluded.} \\
\\ 
\texttt{**Recap of the Setup:**} \\
\texttt{*   **Items Available:** 2 shells, 2 spoons, 1 envelope, 1 pepper.} \\
\texttt{*   **My Value Function:** Shell: 3, Spoon: 0, Envelope: 0, Pepper: 4.} \\
\texttt{*   **Your Value Function (Inferred from your previous calculation):** Shell: 2, Spoon: 2, Envelope: 4, Pepper: 0. (Note: The system's previous output mentioned "shell: 1" in a correction note, but your calculation used 2. Let's stick to the logic that maximized the total to 14).} \\
\\ 
\texttt{**Analyzing the "Total 14" Outcome:**} \\
\texttt{To get a total value of 14, we need to sum the values of all items distributed:} \\
\texttt{*   Total Value of all items based on *my* values: \$2(3) + 2(0) + 1(0) + 1(4) = 10\$.} \\
\texttt{*   Total Value of all items based on *your* values (assuming Shell=2, Spoon=2, Envelope=4, Pepper=0): \$2(2) + 2(2) + 1(4) + 1(0) = 4 + 4 + 4 = 12\$.} \\
\texttt{*   Wait, the prompt says "maximize the sum of your score and the score of the other player." This usually implies using *each person's own* valuation for their own items.} \\
\texttt{    *   If I get items worth \$V\_\{me\}\$ and you get items worth \$V\_\{you\}\$, the total score is \$V\_\{me\} + V\_\{you\}\$.} \\
\texttt{    *   To maximize this sum, we should give items to the person who values them most.} \\
\texttt{    *   **Pepper:** I value it 4, you value it 0. {-}> Give to **Me**.} \\
\texttt{    *   **Shells:** I value them 3 each (total 6), you value them 2 each (total 4). {-}> Give to **Me**.} \\
\texttt{    *   **Spoons:** I value them 0, you value them 2 each (total 4). {-}> Give to **You**.} \\
\texttt{    *   **Envelope:** I value it 0, you value it 4. {-}> Give to **You**.} \\
\\ 
\texttt{**Optimal Allocation for Max Total Score:**} \\
\texttt{*   **Me:** 1 Pepper (4) + 2 Shells (6) = **10**.} \\
            }
        }
    }
    & & \\ \\

    \theutterance \stepcounter{utterance}  
    & & \multicolumn{4}{p{0.6\linewidth}}{
        \cellcolor[rgb]{0.9,0.9,0.9}{
            \makecell[{{p{\linewidth}}}]{
                \texttt{\tiny{[P1$\langle$GM]}}
                \texttt{There seems to be some confusion in the previous turn's breakdown regarding the specific items allocated and the value functions applied. Let's clarify the situation based on the rules and the initial setup to ensure we are on the same page for any future interactions, though the game appears to have concluded.} \\
\\ 
\texttt{**Recap of the Setup:**} \\
\texttt{*   **Items Available:** 2 shells, 2 spoons, 1 envelope, 1 pepper.} \\
\texttt{*   **My Value Function:** Shell: 3, Spoon: 0, Envelope: 0, Pepper: 4.} \\
\texttt{*   **Your Value Function (Inferred from your previous calculation):** Shell: 2, Spoon: 2, Envelope: 4, Pepper: 0. (Note: The system's previous output mentioned "shell: 1" in a correction note, but your calculation used 2. Let's stick to the logic that maximized the total to 14).} \\
\\ 
\texttt{**Analyzing the "Total 14" Outcome:**} \\
\texttt{To get a total value of 14, we need to sum the values of all items distributed:} \\
\texttt{*   Total Value of all items based on *my* values: \$2(3) + 2(0) + 1(0) + 1(4) = 10\$.} \\
\texttt{*   Total Value of all items based on *your* values (assuming Shell=2, Spoon=2, Envelope=4, Pepper=0): \$2(2) + 2(2) + 1(4) + 1(0) = 4 + 4 + 4 = 12\$.} \\
\texttt{*   Wait, the prompt says "maximize the sum of your score and the score of the other player." This usually implies using *each person's own* valuation for their own items.} \\
\texttt{    *   If I get items worth \$V\_\{me\}\$ and you get items worth \$V\_\{you\}\$, the total score is \$V\_\{me\} + V\_\{you\}\$.} \\
\texttt{    *   To maximize this sum, we should give items to the person who values them most.} \\
\texttt{    *   **Pepper:** I value it 4, you value it 0. {-}> Give to **Me**.} \\
\texttt{    *   **Shells:** I value them 3 each (total 6), you value them 2 each (total 4). {-}> Give to **Me**.} \\
\texttt{    *   **Spoons:** I value them 0, you value them 2 each (total 4). {-}> Give to **You**.} \\
\texttt{    *   **Envelope:** I value it 0, you value it 4. {-}> Give to **You**.} \\
\\ 
\texttt{**Optimal Allocation for Max Total Score:**} \\
\texttt{*   **Me:** 1 Pepper (4) + 2 Shells (6) = **10**.} \\
            }
        }
    }
    & \\ \\

    \theutterance \stepcounter{utterance}  
    & \multicolumn{4}{p{0.6\linewidth}}{
        \cellcolor[rgb]{0.9,0.9,0.9}{
            \makecell[{{p{\linewidth}}}]{
                \texttt{\tiny{[P1$\rangle$GM]}}
                \texttt{You are absolutely right to point out the confusion. The previous turn's breakdown contained significant errors regarding the value functions and the final allocation, which led to an incorrect calculation of the total score.} \\
\\ 
\texttt{Let's re{-}evaluate the situation strictly based on the rules and the optimal strategy for maximizing the **sum of our scores** (\$V\_\{me\} + V\_\{you\}\$).} \\
\\ 
\texttt{\#\#\# 1. Clarifying Value Functions} \\
\texttt{Based on your correction and the logic required to reach a total of 14 (or higher), here are the correct valuations we must use:} \\
\\ 
\texttt{*   **My Values:**} \\
\texttt{    *   Shell: 3} \\
\texttt{    *   Spoon: 0} \\
\texttt{    *   Envelope: 0} \\
\texttt{    *   Pepper: 4} \\
\texttt{*   **Your Values:**} \\
\texttt{    *   Shell: 2} \\
\texttt{    *   Spoon: 2} \\
\texttt{    *   Envelope: 4} \\
\texttt{    *   Pepper: 0} \\
\\ 
\texttt{\#\#\# 2. Calculating the Optimal Allocation} \\
\texttt{To maximize the sum of our scores, every item should go to the person who values it most highly.} \\
\\ 
\texttt{*   **Pepper (1 available):**} \\
\texttt{    *   My value: 4} \\
\texttt{    *   Your value: 0} \\
\texttt{    *   **Winner:** Me (+4 to total)} \\
\texttt{*   **Shells (2 available):**} \\
\texttt{    *   My value: 3 each (Total 6)} \\
\texttt{    *   Your value: 2 each (Total 4)} \\
\texttt{    *   **Winner:** Me (+6 to total)} \\
\texttt{*   **Spoons (2 available):**} \\
\texttt{    *   My value: 0 each (Total 0)} \\
\texttt{    *   Your value: 2 each (Total 4)} \\
\texttt{    *   **Winner:** You (+4 to total)} \\
\texttt{*   **Envelope (1 available):**} \\
\texttt{    *   My value: 0} \\
\texttt{    *   Your value: 4} \\
\texttt{    *   **Winner:** You (+4 to total)} \\
\\ 
\texttt{\#\#\# 3. Final Score Calculation} \\
\texttt{*   **My Total Score:** 1 Pepper (4) + 2 Shells (6) = **10**} \\
\texttt{*   **Your Total Score:** 2 Spoons (4) + 1 Envelope (4) = **8**} \\
\texttt{*   **Grand Total:** 10 + 8 = **18**} \\
\\ 
\texttt{*(Note: The previous "Total 14" mentioned in the game log was indeed incorrect based on these specific value functions. The maximum possible sum is 18.)*} \\
\\ 
\texttt{\#\#\# 4. The Correct Proposal} \\
\texttt{Since the game allows us to make a secret proposal at any time to finalize the deal, and we have identified the optimal split, I should have submitted a secret proposal immediately after realizing the error in the previous turn.} \\
\\ 
\texttt{**The Secret Proposal I should have made:**} \\
\texttt{`{[}Proposal:} \\
            }
        }
    }
    & & \\ \\

    \theutterance \stepcounter{utterance}  
    & & & \multicolumn{2}{p{0.3\linewidth}}{
        \cellcolor[rgb]{0.9,0.9,0.9}{
            \makecell[{{p{\linewidth}}}]{
                \texttt{\tiny{[GM$|$GM]}}
                \texttt{continue} \\
            }
        }
    }
    & & \\ \\

    \theutterance \stepcounter{utterance}  
    & & & \multicolumn{2}{p{0.3\linewidth}}{
        \cellcolor[rgb]{0.9,0.9,0.9}{
            \makecell[{{p{\linewidth}}}]{
                \texttt{\tiny{[GM$|$GM]}}
                \texttt{You are absolutely right to point out the confusion. The previous turn's breakdown contained significant errors regarding the value functions and the final allocation, which led to an incorrect calculation of the total score.} \\
\\ 
\texttt{Let's re{-}evaluate the situation strictly based on the rules and the optimal strategy for maximizing the **sum of our scores** (\$V\_\{me\} + V\_\{you\}\$).} \\
\\ 
\texttt{\#\#\# 1. Clarifying Value Functions} \\
\texttt{Based on your correction and the logic required to reach a total of 14 (or higher), here are the correct valuations we must use:} \\
\\ 
\texttt{*   **My Values:**} \\
\texttt{    *   Shell: 3} \\
\texttt{    *   Spoon: 0} \\
\texttt{    *   Envelope: 0} \\
\texttt{    *   Pepper: 4} \\
\texttt{*   **Your Values:**} \\
\texttt{    *   Shell: 2} \\
\texttt{    *   Spoon: 2} \\
\texttt{    *   Envelope: 4} \\
\texttt{    *   Pepper: 0} \\
\\ 
\texttt{\#\#\# 2. Calculating the Optimal Allocation} \\
\texttt{To maximize the sum of our scores, every item should go to the person who values it most highly.} \\
\\ 
\texttt{*   **Pepper (1 available):**} \\
\texttt{    *   My value: 4} \\
\texttt{    *   Your value: 0} \\
\texttt{    *   **Winner:** Me (+4 to total)} \\
\texttt{*   **Shells (2 available):**} \\
\texttt{    *   My value: 3 each (Total 6)} \\
\texttt{    *   Your value: 2 each (Total 4)} \\
\texttt{    *   **Winner:** Me (+6 to total)} \\
\texttt{*   **Spoons (2 available):**} \\
\texttt{    *   My value: 0 each (Total 0)} \\
\texttt{    *   Your value: 2 each (Total 4)} \\
\texttt{    *   **Winner:** You (+4 to total)} \\
\texttt{*   **Envelope (1 available):**} \\
\texttt{    *   My value: 0} \\
\texttt{    *   Your value: 4} \\
\texttt{    *   **Winner:** You (+4 to total)} \\
\\ 
\texttt{\#\#\# 3. Final Score Calculation} \\
\texttt{*   **My Total Score:** 1 Pepper (4) + 2 Shells (6) = **10**} \\
\texttt{*   **Your Total Score:** 2 Spoons (4) + 1 Envelope (4) = **8**} \\
\texttt{*   **Grand Total:** 10 + 8 = **18**} \\
\\ 
\texttt{*(Note: The previous "Total 14" mentioned in the game log was indeed incorrect based on these specific value functions. The maximum possible sum is 18.)*} \\
\\ 
\texttt{\#\#\# 4. The Correct Proposal} \\
\texttt{Since the game allows us to make a secret proposal at any time to finalize the deal, and we have identified the optimal split, I should have submitted a secret proposal immediately after realizing the error in the previous turn.} \\
\\ 
\texttt{**The Secret Proposal I should have made:**} \\
\texttt{`{[}Proposal:} \\
            }
        }
    }
    & & \\ \\

    \theutterance \stepcounter{utterance}  
    & & \multicolumn{4}{p{0.6\linewidth}}{
        \cellcolor[rgb]{0.9,0.9,0.9}{
            \makecell[{{p{\linewidth}}}]{
                \texttt{\tiny{[GM$\rangle$P2]}}
                \texttt{You are absolutely right to point out the confusion. The previous turn's breakdown contained significant errors regarding the value functions and the final allocation, which led to an incorrect calculation of the total score.} \\
\\ 
\texttt{Let's re{-}evaluate the situation strictly based on the rules and the optimal strategy for maximizing the **sum of our scores** (\$V\_\{me\} + V\_\{you\}\$).} \\
\\ 
\texttt{\#\#\# 1. Clarifying Value Functions} \\
\texttt{Based on your correction and the logic required to reach a total of 14 (or higher), here are the correct valuations we must use:} \\
\\ 
\texttt{*   **My Values:**} \\
\texttt{    *   Shell: 3} \\
\texttt{    *   Spoon: 0} \\
\texttt{    *   Envelope: 0} \\
\texttt{    *   Pepper: 4} \\
\texttt{*   **Your Values:**} \\
\texttt{    *   Shell: 2} \\
\texttt{    *   Spoon: 2} \\
\texttt{    *   Envelope: 4} \\
\texttt{    *   Pepper: 0} \\
\\ 
\texttt{\#\#\# 2. Calculating the Optimal Allocation} \\
\texttt{To maximize the sum of our scores, every item should go to the person who values it most highly.} \\
\\ 
\texttt{*   **Pepper (1 available):**} \\
\texttt{    *   My value: 4} \\
\texttt{    *   Your value: 0} \\
\texttt{    *   **Winner:** Me (+4 to total)} \\
\texttt{*   **Shells (2 available):**} \\
\texttt{    *   My value: 3 each (Total 6)} \\
\texttt{    *   Your value: 2 each (Total 4)} \\
\texttt{    *   **Winner:** Me (+6 to total)} \\
\texttt{*   **Spoons (2 available):**} \\
\texttt{    *   My value: 0 each (Total 0)} \\
\texttt{    *   Your value: 2 each (Total 4)} \\
\texttt{    *   **Winner:** You (+4 to total)} \\
\texttt{*   **Envelope (1 available):**} \\
\texttt{    *   My value: 0} \\
\texttt{    *   Your value: 4} \\
\texttt{    *   **Winner:** You (+4 to total)} \\
\\ 
\texttt{\#\#\# 3. Final Score Calculation} \\
\texttt{*   **My Total Score:** 1 Pepper (4) + 2 Shells (6) = **10**} \\
\texttt{*   **Your Total Score:** 2 Spoons (4) + 1 Envelope (4) = **8**} \\
\texttt{*   **Grand Total:** 10 + 8 = **18**} \\
\\ 
\texttt{*(Note: The previous "Total 14" mentioned in the game log was indeed incorrect based on these specific value functions. The maximum possible sum is 18.)*} \\
\\ 
\texttt{\#\#\# 4. The Correct Proposal} \\
\texttt{Since the game allows us to make a secret proposal at any time to finalize the deal, and we have identified the optimal split, I should have submitted a secret proposal immediately after realizing the error in the previous turn.} \\
\\ 
\texttt{**The Secret Proposal I should have made:**} \\
\texttt{`{[}Proposal:} \\
            }
        }
    }
    & \\ \\

    \theutterance \stepcounter{utterance}  
    & & & \multicolumn{4}{p{0.6\linewidth}}{
        \cellcolor[rgb]{0.9,0.9,0.9}{
            \makecell[{{p{\linewidth}}}]{
                \texttt{\tiny{[GM$\langle$P2]}}
                \texttt{{[}Proposal: 1 pepper, 2 shells{]}} \\
            }
        }
    }
     \\ \\

    \theutterance \stepcounter{utterance}  
    & & & \multicolumn{2}{p{0.3\linewidth}}{
        \cellcolor[rgb]{0.9,0.9,0.9}{
            \makecell[{{p{\linewidth}}}]{
                \texttt{\tiny{[GM$|$GM]}}
                \texttt{proposal} \\
            }
        }
    }
    & & \\ \\

    \theutterance \stepcounter{utterance}  
    & & & \multicolumn{2}{p{0.3\linewidth}}{
        \cellcolor[rgb]{0.9,0.9,0.9}{
            \makecell[{{p{\linewidth}}}]{
                \texttt{\tiny{[GM$|$GM]}}
                [2, 0, 0, 1]
            }
        }
    }
    & & \\ \\

    \theutterance \stepcounter{utterance}  
    & & \multicolumn{4}{p{0.6\linewidth}}{
        \cellcolor[rgb]{0.9,0.9,0.9}{
            \makecell[{{p{\linewidth}}}]{
                \texttt{\tiny{[P1$\langle$GM]}}
                \texttt{The other player has submitted a secret proposal. It is now your turn to submit a proposal of your own.} \\
            }
        }
    }
    & \\ \\

    \theutterance \stepcounter{utterance}  
    & \multicolumn{4}{p{0.6\linewidth}}{
        \cellcolor[rgb]{0.9,0.9,0.9}{
            \makecell[{{p{\linewidth}}}]{
                \texttt{\tiny{[P1$\rangle$GM]}}
                \texttt{{[}Proposal: 1 pepper, 2 shells{]}} \\
            }
        }
    }
    & & \\ \\

    \theutterance \stepcounter{utterance}  
    & & & \multicolumn{2}{p{0.3\linewidth}}{
        \cellcolor[rgb]{0.9,0.9,0.9}{
            \makecell[{{p{\linewidth}}}]{
                \texttt{\tiny{[GM$|$GM]}}
                \texttt{proposal} \\
            }
        }
    }
    & & \\ \\

    \theutterance \stepcounter{utterance}  
    & & & \multicolumn{2}{p{0.3\linewidth}}{
        \cellcolor[rgb]{0.9,0.9,0.9}{
            \makecell[{{p{\linewidth}}}]{
                \texttt{\tiny{[GM$|$GM]}}
                [2, 0, 0, 1]
            }
        }
    }
    & & \\ \\

    \theutterance \stepcounter{utterance}  
    & & & \multicolumn{2}{p{0.3\linewidth}}{
        \cellcolor[rgb]{0.9,0.9,0.9}{
            \makecell[{{p{\linewidth}}}]{
                \texttt{\tiny{[GM$|$GM]}}
                [[2, 0, 0, 1], [2, 0, 0, 1]]
            }
        }
    }
    & & \\ \\

    \theutterance \stepcounter{utterance}  
    & \multicolumn{4}{p{0.6\linewidth}}{
        \cellcolor[rgb]{0.95,0.95,0.95}{
            \makecell[{{p{\linewidth}}}]{
                \texttt{\tiny{[P2$\rightarrow$P1]}}
                \texttt{Post{-}episode Reflexion (not sent to GM):} \\
\texttt{New plan: I failed because I submitted a secret proposal (`{[}Proposal: 1 pepper, 2 shells{]}`) that was incomplete and did not specify the allocation for the other player (the Game Master). The rules state that proposals must be complementary to succeed; by only listing items for myself, the system could not verify if the remaining items (2 spoons, 1 envelope) were sufficient and correctly assigned to the other player to form a valid, complementary deal. To succeed, I must ensure my secret proposal explicitly defines the full split: clearly stating which items I want and which items the other player gets, ensuring the sum of requested items equals the total available inventory (2 shells, 2 spoons, 1 envelope, 1 pepper). I will draft a proposal like `{[}Proposal: 1 pepper, 2 shells for me; 2 spoons, 1 envelope for you{]}` to guarantee complementarity.} \\
            }
        }
    }
    & & \\ \\

    \theutterance \stepcounter{utterance}  
    & & & \multicolumn{4}{p{0.6\linewidth}}{
        \cellcolor[rgb]{0.95,0.95,0.95}{
            \makecell[{{p{\linewidth}}}]{
                \texttt{\tiny{[P1$\rightarrow$P2]}}
                \texttt{Post{-}episode Reflexion (not sent to GM):} \\
\texttt{New plan: I failed because I spent my turns analyzing the math and explaining the error instead of immediately submitting the required secret proposal format. The game rules state that to finalize the deal, I must use the exact string `{[}Proposal: <items>{]}`. In my last turn, I wrote the proposal but did not include the opening square bracket `{[}`, causing the system to reject it as a valid secret proposal. Next time, upon identifying the optimal allocation (Me: 1 pepper, 2 shells; You: 2 spoons, 1 envelope), I will immediately output the exact string `{[}Proposal: 1 pepper, 2 shells{]}` without any preceding text or analysis to ensure the game terminates successfully with the maximum combined score.} \\
            }
        }
    }
     \\ \\

\end{supertabular}
}

\end{document}
